\paragraph{An assessment closes the interval it measures.}
A tax assessment differs from a payment because it measures activity over an
interval and constrains which events may follow. The optional
\texttt{MixedEventTimeline/2} retains the same native annual book and decision
rules while adding an explicit clock for current assessment. The interval
has an inclusive start, an exclusive end and a final included activity date.
That final date must be the day before the exclusive end, and it must equal
the date supplied in the timeline. No extra \code{event.date} is introduced.
The assessment reads the book at its actual turn, after earlier successful
actions; static validation does not calculate tax against the initial book.

Consider the same opening cash of USD~452, inventory of USD~80 and carried tax
payable of USD~10.25, now with an explicitly enabled one-unit sale on
2~January 2027. The sale earns USD~15 against cost of USD~10. Under the
constructed assumption of zero tax adjustments and a one-quarter tax rate,
a subsequent assessment through that day measures profit of USD~5 and current
tax of USD~1.25. Cash is USD~467, inventory is USD~70 and total tax payable is
USD~11.50. Assessment creates no cash flow and preserves the carried vintages.
These numbers illustrate the supplied tax model; they are not an assessment
of a named company's liability.

Reversing the explicit same-day priorities gives a different admissibility
result. Assessment first records zero current profit and closes the interval
through 2~January. The later sale is refused because it would add activity
inside an already assessed interval. The runner preserves the successful
assessment, cash of USD~452 and the exact rejected sale. It neither silently
moves the sale ahead of assessment nor rolls back the valid prefix.

A half-unit sale on the next day, after the successful sale-first sequence,
adds proceeds of USD~7.50 and cost of USD~5. Cumulative profit becomes USD~7.50
and cumulative tax USD~1.875. Replacing the earlier assessment therefore
records only USD~0.625 of additional expense and updates the same current
vintage. Cash becomes USD~474.50 and inventory USD~65. This replacement rule
prevents the earlier USD~1.25 charge from being counted again.

Twelve focused tests and 49 regression tests preserve eight complete exchanges
in \path{data/accounting-master/assessment-timeline-20261025/}. The comparison
checks the full materialized decision, native transition, journals and report,
alongside independent arithmetic. Unknown rates retain their unresolved
status after a successful sale; malformed clocks, stale initial state and
ambiguous priorities refuse. The existing bounds of 32 native entries and
64 total actions are tested with actual boundary cases. Deferred measurement
still requires its own state-binding rule, and carry still requires a fiscal
boundary rule. Completing this declared timeline therefore establishes neither
complete fiscal activity nor broader accounting admission.
